\chapter{Burning Feet}

\section*{Definition}
A sensation of heat, burning, or warmth in the feet, often accompanied by tingling, numbness, or sharp pain, indicating underlying peripheral nerve dysfunction.

\section*{Pathophysiology}
Burning feet most commonly result from peripheral neuropathy, particularly length-dependent axonal neuropathy where distal nerve endings are affected first. Metabolic (diabetes), toxic (alcohol), nutritional (B vitamin deficiencies), and inflammatory (autoimmune) causes damage to small nerve fibers (A-delta and C fibers) that transmit pain and temperature sensation, leading to aberrant nociceptive signaling.

\section*{Etiology}
\begin{itemize}
  \item Diabetic peripheral neuropathy (most common)
  \item Alcoholic neuropathy
  \item Vitamin B deficiency (B1, B6, B12)
  \item Peripheral arterial disease (critical limb ischemia)
  \item Tarsal tunnel syndrome
  \item Hypothyroidism
  \item Chemotherapy-induced peripheral neuropathy
  \item Chronic kidney disease (uremic neuropathy)
  \item Erythromelalgia (rare vascular disorder)
  \item Heavy metal toxicity (lead, mercury, arsenic)
\end{itemize}

\section*{Indications for Evaluation}
Seek care if:
\begin{itemize}
  \item Severe or worsening symptoms
  \item Burning feet with uncontrolled diabetes, progressive weakness, skin breakdown, or weight loss
  \item Accompanied by other concerning signs
\end{itemize}

\section*{Related Symptoms}
\begin{itemize}
  \item Foot numbness
  \item Tingling
  \item Leg pain
  \item Skin changes
  \item Foot ulcers
\end{itemize}

\textbf{Note:} Always consider serious underlying conditions when evaluating persistent burning feet.

\section*{Self-Care / Home Management}
\begin{itemize}
  \item Rest and avoid activities that may aggravate the condition.
  \item Maintain adequate hydration and a balanced diet.
  \item Over-the-counter remedies may provide symptomatic relief (consult a pharmacist or doctor).
  \item Monitor symptoms and seek professional advice if they persist or worsen.
  \item Avoid self-medicating with prescription medications without medical guidance.
\end{itemize}

\section*{Diagnostic Approach}
\begin{itemize}
  \item A thorough history and physical examination are the first steps in evaluation.
  \item Laboratory tests (blood work, urinalysis) may be ordered based on clinical suspicion.
  \item Imaging studies (X-ray, ultrasound, CT, MRI) may be indicated when structural pathology is suspected.
  \item Specialist referral is arranged when the diagnosis remains uncertain or requires specific expertise.
\end{itemize}

\section*{Treatment / Management Overview}
\begin{itemize}
  \item Treatment targets the underlying cause identified through diagnostic evaluation.
  \item Supportive care and symptom management are provided while awaiting definitive therapy.
  \item Medications, lifestyle modifications, and physical therapies may be prescribed as appropriate.
  \item Follow-up ensures treatment effectiveness and allows timely adjustments to the care plan.
\end{itemize}

\clearpage